\section{Contribuições do trabalho}

Este trabalho apresenta contribuições relacionadas à modelagem, simulação e análise de uma missão de \gls{rvdb} envolvendo uma espaçonave equipada com um \gls{mr}. 
As principais contribuições podem ser resumidas nos seguintes pontos:

\begin{itemize}

\item Desenvolvimento de um ambiente de simulação capaz de representar as diferentes fases de uma missão de \gls{rvdb}, incluindo aproximação orbital, aproximação relativa e operação do \gls{mr}.

\item Integração, em um único modelo de simulação, da dinâmica orbital relativa baseada nas equações de \gls{hcw}, da dinâmica de atitude da espaçonave e da dinâmica acoplada entre a espaçonave e o \gls{mr}.

\item Implementação de leis de controle \gls{pd} para o controle translacional relativo e para o controle de atitude da espaçonave, permitindo avaliar o comportamento do sistema durante as fases críticas da missão.

\item Modelagem do efeito de torques de reação gerados pelo movimento do \gls{mr} sobre a dinâmica de atitude da espaçonave, incluindo sua compensação no sistema de controle.

\item Estruturação do modelo de simulação no ambiente Simulink de forma modular, permitindo a separação entre os diferentes modos de operação do sistema e a comutação entre diferentes modelos dinâmicos ao longo da missão.

\end{itemize}

Essas contribuições permitem analisar de forma integrada os efeitos da interação entre dinâmica orbital, dinâmica de atitude e movimentação do \gls{mr} em missões de \gls{rvdb}.